% Linha do tempo do data warehouse (cap. 2). Fontes na legenda do capitulo.
\begin{tikzpicture}[diagrama, x=1mm, y=1mm]
  \draw[fio, line width=0.8pt] (18,5) -- (18,-85);
  \foreach \y/\ano/\texto/\peso in {
      0/1958/{Luhn descreve um \enquote{sistema de business intelligence} na IBM}/normal,
    -10/1970/{Codd propõe o modelo relacional, base dos bancos de dados SQL}/normal,
    -20/1988/{Devlin e Murphy descrevem a arquitetura de um \enquote{business data warehouse} na IBM}/normal,
    -30/1992/{Inmon publica \emph{Building the Data Warehouse}: o armazém corporativo, normalizado, construído de cima para baixo}/destaque,
    -40/1996/{Kimball publica \emph{The Data Warehouse Toolkit}: a modelagem dimensional, construída de baixo para cima}/destaque,
    -50/{Anos 2000}/{O Data Vault, de Dan Linstedt, ganha espaço como alternativa para a camada de integração}/normal,
    -60/{Anos 2010}/{Armazéns de dados na nuvem separam armazenamento de processamento, e o ETL dá lugar ao ELT}/normal,
    -70/2021/{Armbrust e colegas propõem o lakehouse, que junta data lake e data warehouse}/normal,
    -80/{Anos 2020}/{A IA generativa entra no fluxo de trabalho com dados}/normal%
  }{
    \ifthenelse{\equal{\peso}{destaque}}
      {\filldraw[draw=tinta, fill=tinta] (18,\y) circle (1.4);}
      {\filldraw[draw=cinza, fill=white, line width=0.6pt] (18,\y) circle (1.1);}
    \node[anchor=east, font=\footnotesize] at (14.5,\y) {\ano};
    \node[anchor=west, text width=92mm, align=left, font=\footnotesize] at (22,\y) {\texto};
  }
\end{tikzpicture}
